\section{Testing with one clean qubit}\label{sec:clean}
Recent work of Shao \cite{shao} studies quantum query algorithms with one initially pure qubit, any number of maximally mixed qubits, input-independent unitaries, controlled phase queries, and a final measurement. We use this model with unbounded mixed workspace. Its queries can be interleaved with arbitrary computation. Limiting the initial purity changes the answer to our testing problem.
\begin{theorem}
In the range of Theorem~\ref{th:nonadaptive}, the bounded-error one-clean-qubit query cost of testing Boolean MCSP is
\[
                  \Theta_{\eta,\epsilon,\theta}\bigl(\min\{N,T/\xi\}\bigr).
\]
The result holds for the exact yes side and for the tolerant yes side specified there.
\end{theorem}

\paragraph{A lower bound for inputs supported on a block.}
Let $g$ be a partial Boolean function on $M$ bits. In an oracle of length $N$, suppose a known set $B$ contains copies of the $M$ bits according to a known map and all positions outside $B$ are fixed. Put $p=|B|/N$. If $g$ has ordinary bounded-error quantum query complexity $d$, evaluating this embedded problem with $k\ge1$ initially pure qubits and all remaining qubits maximally mixed takes
\[
                    \Omega\!\left(\frac{d}{\min\{1,2^k p\}}\right)
                    \quad\text{queries}.                                      \tag{3}
\]
The final measurement can be arbitrary. To see the bound, use the general adversary characterization for partial Boolean functions \cite{adv}. Take a Hermitian adversary matrix $\Gamma$, zero on equal output labels, and normalize $\max_j\|\Gamma\circ\Delta_j\|=1$, where $\Delta_j(x,y)=\mathbf1[x_j\ne y_j]$. We can arrange $\|\Gamma\|=\Omega(d)$ and take a unit eigenvector $v$ for a positive largest eigenvalue, changing the sign if needed.

Purify the same mixed initial state for every input. The ordinary adversary progress is the average Gram inner product of the resulting states weighted by $\overline v_x\Gamma_{xy}v_y$. Its initial value is $\|\Gamma\|$. A successful final measurement changes it by $\Omega(\|\Gamma\|)$ by the usual bounded-error adversary argument \cite{adv}. Write $\psi_{x,i}$ for the branch of a prequery state at address $i$ with phase control active. A controlled phase query changes progress in absolute value by at most
\[
       2\sum_{i\in B}\sum_x |v_x|^2\|\psi_{x,i}\|^2 .
\]
Indeed, at each $i$ the signed quadratic form has matrix $\Gamma\circ\Delta_{j(i)}$ tensored with the identity on the workspace, of norm at most one. Fixed addresses make zero contribution. On a computational space of dimension $D$, the state on any particular input still has operator norm $2^k/D$ after all preceding unitaries. The projector onto addresses in $B$ has rank $pD$, even if we ignore the phase-control condition. Its mass is at most $\min(1,2^kp)$ for each $x$. Telescoping proves (3).

Use the hard base promise from Theorem~\ref{th:nonadaptive}: a codeword versus a word $\alpha M$-far from every size-$s$ table. Its ordinary $d=\Omega(M)$ follows because the easy law has $\Omega(M)$-wise uniform bits and a uniformly random word is far with probability $1-o(1)$. No polynomial of smaller degree can distinguish these unconditional laws by a constant. For (3) the active prefix may even be fixed in advance. Taking $k=1$ gives $\Omega(M/p)$, as required. The dense case gives $\Omega(N)$ already in the ordinary model.

\paragraph{The matching upper bound.}
Shao's Lemma~3.8 computes an arbitrary Boolean predicate of $q$ bits sampled at independent uniform addresses with $2q$ queries and one clean qubit \cite{shao}. The postprocessing predicate can also depend on the addresses. In the proof, all address registers are preserved on each computational branch, so the input-independent reversible postprocessing can select its predicate according to that branch. This implements the sample tester in Theorem~\ref{th:nonadaptive}. To read the entire table, start with one maximally mixed address $J$. For each $i\in\{0,1\}^n$, temporarily replace $J$ by $J\oplus i$, query it into the corresponding mixed mask bit, and undo the address change. Perform this loading and unloading in the same two places in Shao's construction. Conditional on $J$, the $N$ loaded bits give the whole input in a known order. Use the acceptance predicate in that order. This gives $2N$ queries. The workspace and intervening computation may be very large.
